\documentclass[11pt,a4paper]{article}

\usepackage[T1]{fontenc}
\usepackage[utf8]{inputenc}
\usepackage{lmodern}
\usepackage[ngerman]{babel}
\usepackage[a4paper,margin=2.2cm]{geometry}
\usepackage{array}
\usepackage{longtable}
\usepackage{tabularx}
\usepackage{xcolor}
\usepackage{hyperref}

\setlength{\parindent}{0pt}
\setlength{\parskip}{0.5em}
\renewcommand{\arraystretch}{1.35}

\newcommand{\answerline}[1]{\rule{#1}{0.4pt}}
\newcommand{\term}[1]{\fcolorbox{black!20}{blue!6}{\strut #1}}

\title{\textbf{Aufgabenblatt: Softwaresysteme migrieren}}
\author{Modul 321 -- Lektion 05}
\date{}

\begin{document}

\maketitle

\textbf{Name:} \answerline{7cm} \hfill \textbf{Datum:} \answerline{3.5cm}

\vspace{0.8em}

\fcolorbox{blue!40}{blue!6}{
  \parbox{\dimexpr\linewidth-2\fboxsep-2\fboxrule-6pt\relax}{
    \textbf{Ziel der Lektion:} Sie können zentrale Vorgehensweisen zur Migration
    von Softwaresystemen unterscheiden, Risiken benennen und für einfache
    Szenarien eine sinnvolle Strategie auswählen und begründen.
  }
}

\section*{Hinweise}
\begin{itemize}
  \item Antworten Sie in ganzen Sätzen oder präzisen Stichworten.
  \item Begründen Sie Ihre Entscheidungen knapp, aber nachvollziehbar.
  \item Wo sinnvoll, dürfen Sie kleine Skizzen oder Ablaufdiagramme ergänzen.
\end{itemize}

\section*{1. Begriffe erklären}
Erklären Sie die folgenden Begriffe in \textbf{1 bis 3 eigenen Sätzen}.

\begin{longtable}{|>{\raggedright\arraybackslash}p{0.28\linewidth}|>{\raggedright\arraybackslash}p{0.64\linewidth}|}
\hline
\textbf{Begriff} & \textbf{Ihre Erklärung} \\
\hline
\term{Migration} & \\[2.2cm] \hline
\term{Big Bang} & \\[2.2cm] \hline
\term{Strangler Pattern} & \\[2.2cm] \hline
\term{Parallelbetrieb} & \\[2.2cm] \hline
\end{longtable}

\section*{2. Aussagen zuordnen}
Ordnen Sie jede Aussage der passenden Strategie zu:

\term{Big Bang}, \term{Strangler Pattern}, \term{beide}, \term{keine von beiden}

\begin{enumerate}
  \item Das alte System bleibt zunächst teilweise in Betrieb.
  \item Der Wechsel auf das neue System erfolgt in einem grossen Umschaltmoment.
  \item Übergangsschnittstellen und Routing spielen oft eine wichtige Rolle.
  \item Das Risiko verteilt sich eher auf mehrere kleinere Schritte.
  \item Ein Wartungsfenster am Umstellungstag ist oft besonders wichtig.
  \item Die Strategie kann nur bei Microservices eingesetzt werden.
\end{enumerate}

\section*{3. Big Bang beurteilen}
Beantworten Sie die Fragen kurz.

\begin{enumerate}
  \item Beschreiben Sie den Grundgedanken eines \term{Big Bang}.
  \item Nennen Sie \textbf{zwei Vorteile}.
  \item Nennen Sie \textbf{drei Risiken}.
  \item Für welche Art von System wäre ein Big Bang eher vertretbar?
\end{enumerate}

\section*{4. Strangler Pattern verstehen}
Ein bestehender Webshop soll modernisiert werden. Die Suche und der Produktkatalog
sollen zuerst erneuert werden, der Checkout bleibt vorerst im alten System.

\begin{enumerate}
  \item Warum passt dieses Vorgehen zum \term{Strangler Pattern}?
  \item Was ist der Vorteil, wenn man mit Katalog und Suche beginnt?
  \item Nennen Sie eine technische Herausforderung beim Parallelbetrieb.
  \item Was müsste geschehen, damit der alte Teil später entfernt werden kann?
\end{enumerate}

\section*{5. Vergleichstabelle}
Vergleichen Sie \term{Big Bang} und \term{Strangler Pattern}.

\begin{tabularx}{\linewidth}{|>{\raggedright\arraybackslash}p{0.28\linewidth}|>{\raggedright\arraybackslash}X|>{\raggedright\arraybackslash}X|}
\hline
\textbf{Aspekt} & \textbf{Big Bang} & \textbf{Strangler Pattern} \\
\hline
Zeitpunkt der Umstellung & & \\[1.1cm] \hline
Risiko im produktiven Betrieb & & \\[1.1cm] \hline
Komplexität der Übergangsphase & & \\[1.1cm] \hline
Geeignet für welche Situationen? & & \\[1.1cm] \hline
\end{tabularx}

\section*{6. Szenario entscheiden}
Ein kleines internes Schulverwaltungstool soll ersetzt werden.

Rahmenbedingungen:
\begin{itemize}
  \item nur 200 Benutzerinnen und Benutzer
  \item wenige Funktionen
  \item die Anwendung darf an einem Wochenende kurz offline sein
  \item die Datenstruktur ist überschaubar
\end{itemize}

\begin{enumerate}
  \item Welche Migrationsstrategie ist hier naheliegend?
  \item Begründen Sie die Wahl mit \textbf{zwei Argumenten}.
  \item Nennen Sie trotzdem \textbf{ein Risiko}, das beachtet werden muss.
\end{enumerate}

\section*{7. Kritisches Szenario}
Ein grosser Online-Shop mit vielen täglichen Bestellungen soll modernisiert werden.
Ein Ausfall während des Checkouts wäre teuer.

\begin{enumerate}
  \item Welche Strategie ist hier eher naheliegend?
  \item Warum wäre ein Big Bang in diesem Fall kritisch?
  \item Welchen Teil des Shops würden Sie zuerst migrieren?
  \item Welchen Teil würden Sie eher spät migrieren?
\end{enumerate}

\section*{8. Eigene Migrationsskizze}
Eine bestehende Lernplattform besteht aus:
\begin{itemize}
  \item Login
  \item Kursübersicht
  \item Dateiupload
  \item Nachrichten
  \item Videolektionen
\end{itemize}

Bearbeiten Sie die Aufgaben:
\begin{enumerate}
  \item Wählen Sie \term{Big Bang} oder \term{Strangler Pattern}.
  \item Begründen Sie Ihre Wahl.
  \item Skizzieren oder beschreiben Sie in 3 bis 5 Schritten, wie die Migration ablaufen könnte.
  \item Nennen Sie einen Bereich mit hohem Risiko und einen Bereich mit tieferem Risiko.
\end{enumerate}

\end{document}
